\chapter{Diseño e Implementación del Sistema}

\section{Descripción del Problema}

El diseño de materiales para baterías de ion-litio enfrenta un cuello de
botella combinatorio: solo en óxidos y fosfatos de metales de transición con
litio, el espacio de composiciones posibles supera las 500\,000 combinaciones,
de las cuales solo una fracción mínima ha sido calculada mediante Teoría del
Funcional de la Densidad (DFT) \citep{urban2016computational}. Un cálculo DFT
individual tarda horas o días \citep{jain2013materialsproject}, por lo que
evaluar el espacio de candidatos exhaustivamente es inviable.

El problema concreto: \emph{dado un conjunto de materiales con propiedades
DFT ya calculadas para cátodos y ánodos de Li-ion (LiFePO\textsubscript{4}
\citep{padhi1997phospho}, óxidos Li-Co-Mn-Ni), ¿cómo entrenar modelos de
aprendizaje automático que aprendan la relación entre composición química y
estabilidad termodinámica, prediciendo en microsegundos para composiciones
nuevas sin requerir DFT por candidato?} La estabilidad de un óxido mixto
depende de interacciones no lineales entre campo de ligandos, radio iónico y
electronegatividad \citep{behler2007generalized}; no es un problema de
fórmula cerrada, lo que justifica recurrir a ML supervisado entrenado sobre
datos DFT reales \citep{ward2016general, meredig2014combinatorial}.

\section{Requisitos y Criterios de Diseño}

\textbf{Funcionales:} validación de datasets CSV; generación de descriptores
composicionales; entrenamiento/comparación de modelos de regresión y
clasificación; predicción por lotes con detección de muestras fuera de
dominio (OOD); \textit{ranking} de candidatos por aplicación (cátodo, ánodo,
electrolito sólido).

\textbf{No funcionales:} reproducibilidad (semilla fija, \textit{hash}
SHA-256 de datasets/artefactos); seguridad (autenticación, RBAC, validación
estricta de archivos); trazabilidad (auditoría de cada predicción); honestidad
científica (nunca calibrar una predicción sin justificarlo; advertir
predicciones fuera de dominio).

\textbf{Restricción de seguridad:} \texttt{PyJWT} exclusivo para tokens, \textit{hashing} de contraseñas con
Argon2, y prohibición de ejecutar contenido subido por el usuario.

\section{Diseño de la Solución}

\subsection{Arquitectura General}

Arquitectura limpia por capas: el dominio no depende de infraestructura, la
aplicación solo depende del dominio. Flujo de una predicción:
\texttt{API (FastAPI) → Caso de uso → Servicio de dominio → Repositorio →
PostgreSQL / artefacto \texttt{.joblib}}. Esto permite probar la lógica
científica sin base de datos real y sustituir el motor de ML sin tocar las
rutas HTTP.

\subsection{Selección de Variables y Datos}

Los datos provienen de Materials Project \citep{jain2013materialsproject},
calculados con el funcional PBE-GGA \citep{perdew1996generalized}. El
dataset demo contiene 154 materiales Li-ion (151 válidos tras validación).

\begin{table}[H]
\centering
\small
\caption{Variables objetivo predichas por el sistema.}
\label{tab:variables}
\begin{tabularx}{\textwidth}{l l X}
\toprule
\textbf{Variable} & \textbf{Unidad} & \textbf{Significado físico} \\
\midrule
\texttt{energy\_above\_hull} & eV/átomo & Estabilidad relativa al casco convexo \\
\texttt{formation\_energy\_per\_atom} & eV/átomo & Energía liberada al formarse \\
\texttt{band\_gap} & eV & Carácter conductor/semiconductor/aislante \\
\texttt{is\_stable} & booleano & EAH $\leq 0.05$ eV/átomo \\
\bottomrule
\end{tabularx}
\end{table}

Cada fórmula se transforma en $\sim$57 descriptores composicionales vía
\texttt{pymatgen} (esquema Magpie \citep{ward2016general}, análogo a
\texttt{matminer} \citep{ward2018matminer}): estadísticos de propiedades
elementales ponderados por fracción atómica, más fracciones de elementos
relevantes para baterías (Li, Fe, Mn, Co, Ni, O, P, S).

\subsection{Decisiones de Ingeniería}

Se eligió \texttt{scikit-learn} sobre redes de grafos (CGCNN
\citep{xie2018cgcnn}, MEGNet \citep{chen2019megnet}) porque éstas requieren
estructura cristalina 3D no disponible para composiciones hipotéticas, y
datasets mucho mayores a 151 materiales (limitación discutida en
\ref{sec:limitaciones}). Se comparan cuatro modelos —\textit{Ridge},
\textit{Random Forest}, \textit{Gradient Boosting} y un clasificador
\textit{Random Forest}— siguiendo la metodología de \textit{benchmarking}
de \citet{dunn2020matbench}, en vez de comprometerse con uno solo. Toda
predicción no calibrada se marca como tal (\texttt{is\_calibrated=False}).

\section{Implementación}

\subsection{Preprocesamiento de Datos}

Cada CSV se valida antes de tocar la base de datos: codificación UTF-8,
tamaño máximo, columna \texttt{formula}, rango físico por propiedad, y
\textit{hash} SHA-256 para duplicados. Las filas rechazadas se registran
con su motivo, nunca se descartan en silencio. Las fórmulas se parsean con
\texttt{pymatgen.Composition}; una fórmula inválida se rechaza antes del
pipeline de descriptores.

\subsection{Entrenamiento del Modelo}

El entrenamiento envuelve el estimador en un \textit{pipeline} con
imputación de medianas y escalado robusto, con semilla fija
(\texttt{FIXED\_RANDOM\_SEED=42}) para reproducibilidad bit-a-bit:

\begin{lstlisting}[language=Python, caption={Pipeline de entrenamiento (adaptado de \texttt{infrastructure/ml/trainers.py})}, label={lst:training}]
def train(self, model_type, task_type, X_train, y_train,
          hyperparams=None, scale=True):
    # Merge model-specific defaults with caller overrides
    defaults = (REGRESSOR_DEFAULTS if task_type == TaskType.REGRESSION
                else CLASSIFIER_DEFAULTS).get(model_type, {})
    params = {**defaults, **(hyperparams or {})}
    estimator = _build_estimator(model_type, task_type, params)
    preprocessor = build_sklearn_pipeline(task_type.value, scale=scale)

    # Imputer + scaler + estimator as a single fitted object
    full_pipeline = Pipeline([("preprocessor", preprocessor),
                               ("model", estimator)])
    full_pipeline.fit(X_train, y_train)
    return full_pipeline
\end{lstlisting}

\subsection{Validaciones y Manejo de Errores}

Todo artefacto \texttt{.joblib} se verifica por \textit{hash} SHA-256 antes
de cargarse en memoria (\texttt{load\_artifact}), evitando ejecutar un
artefacto alterado; un \textit{hash} no coincidente lanza
\texttt{ArtifactIntegrityError} sin cargar el archivo. El sistema define
además una jerarquía de excepciones tipadas (\texttt{InvalidChemicalFormu\-laError},
\texttt{InsufficientDataError}, \texttt{TargetLeakageError}) que separan el
mensaje seguro mostrado al usuario del detalle técnico registrado solo en
los \textit{logs} del servidor, evitando filtrar trazas internas.

\subsection{Módulo de Trabajos DFT (Etapa 13)}

La plataforma incluye además un módulo operativo de orquestación de cálculos
DFT (\texttt{/api/v1/dft-jobs}: envío, listado, estado, cancelación, ingesta
de resultados y descarga de \textit{inputs}), con seguimiento de cada
trabajo en PostgreSQL y una página dedicada en el \textit{frontend}
(\textit{Trabajos DFT}). Un adaptador local ejecuta cálculos reales vía
ASE/EMT para sistemas metálicos (Al, Cu, Ag, Au, Ni, Pd, Pt, Zn, Cd, Hg) y
recurre a una aproximación determinista —marcada explícitamente como no
equivalente a DFT real— para el resto; un adaptador SLURM genera archivos
VASP (POSCAR/INCAR/KPOINTS vía \texttt{pymatgen}) y \textit{scripts}
\texttt{sbatch} reales, sin conexión aún a un clúster HPC. Esta capa está
disponible en la plataforma, pero los resultados de la sección siguiente
provienen exclusivamente de datos DFT precalculados de Materials Project,
no de cálculos ejecutados por este módulo.

\section{Evaluación y Resultados}

\subsection{Métricas de Desempeño}

\begin{table}[H]
\centering
\small
\caption{Desempeño real de los modelos activos sobre el conjunto de prueba (151 materiales, \textit{split} 80/20, semilla fija).}
\label{tab:resultados}
\begin{tabularx}{\textwidth}{l l X}
\toprule
\textbf{Objetivo / Modelo} & \textbf{Métrica de prueba} & \textbf{Valor} \\
\midrule
EAH (Gradient Boosting) & MAE / RMSE / $R^2$ & 0.0186 / 0.0293 eV/at. / 0.155 \\
$E_f$ por átomo (Random Forest) & MAE / RMSE / $R^2$ & 0.210 / 0.301 eV/at. / 0.869 \\
\texttt{is\_stable} (RF, clasificación) & F1-macro / ROC-AUC & 0.583 / 0.692 \\
\bottomrule
\end{tabularx}
\end{table}

\subsection{Análisis de Resultados}

La energía de formación se predice con alta precisión ($R^2=0.869$),
consistente con \citet{meredig2014combinatorial}. \texttt{energy\_above\_hull}
—la propiedad de estabilidad más relevante para ciclado— es más difícil
($R^2=0.155$): es relativa al casco convexo completo del sistema químico, no
una propiedad intrínseca, y con 151 muestras el modelo generaliza pobremente
entre familias químicas. El clasificador binario (F1-macro$=0.583$) es más
robusto que la regresión continua equivalente al colapsar el problema a una
decisión binaria con \texttt{class\_weight="balanced"}. La validación
cruzada (MAE$_{cv}=0.0166$) confirma que estas cifras son estables, no
producto de un \textit{split} afortunado.

\subsection{Limitaciones del Sistema}
\label{sec:limitaciones}

\textbf{(1) Tamaño del dataset:} 151 materiales es insuficiente para
generalizar fuera de las familias observadas (Li-Fe-O, Li-Mn-O, Li-Co-O); el
sistema marca estas predicciones como OOD en vez de reportarlas con falsa
confianza. \textbf{(2) Descriptores composicionales:} dos polimorfos de la
misma fórmula (p.\,ej.\ TiO\textsubscript{2} rutilo vs.\ anatasa) son
indistinguibles sin información estructural 3D —lo que CGCNN
\citep{xie2018cgcnn} resuelve a costa de requerir estructura conocida.
\textbf{(3) Brecha estabilidad–desempeño:} el sistema predice si un material
\emph{existe y es estable} (precondición necesaria), no su voltaje,
capacidad específica o conductividad iónica \citep{urban2016computational,
bachman2016inorganic}, que requieren simulaciones adicionales no incluidas
en el dataset. Además, los datos heredan el error sistemático del funcional
PBE \citep{perdew1996generalized}, y EAH$=0$ no garantiza sintetizabilidad
experimental \citep{sun2016thermodynamic, aykol2018thermodynamic}.

\section{Conclusiones}

El sistema demuestra que un \textit{pipeline} clásico de ML —sin redes
generativas ni de grafos— sustituye un cálculo DFT (microsegundos vs.\
horas) para cribar estabilidad termodinámica de candidatos a cátodos y
ánodos de ion-litio, cumpliendo el objetivo de reducir el costo de explorar
el espacio composicional. El desempeño —alto para energía de formación,
moderado para estabilidad— es honesto dado el tamaño del dataset, y se
comunica explícitamente en vez de ocultarse. La Etapa 13 —el módulo de orquestación de
trabajos DFT descrito en Implementación— ya está operativa; el trabajo
futuro es conectarla a un clúster HPC real, incorporar sus resultados al
dataset de entrenamiento, y sumar descriptores estructurales o redes de
grafos cristalinos (CGCNN, ALIGNN) cuando exista información 3D suficiente,
cerrando progresivamente la brecha hacia la predicción directa de
desempeño de batería y no solo de estabilidad termodinámica.